\documentclass[12pt,a4paper]{article}
\usepackage{amsmath,amssymb,amsfonts}
\usepackage{geometry}
\geometry{margin=2.5cm}
\title{\bf Helical Geometry of the Riemann Zeta Function:\\ 
The Conical Helix and the Hilbert--P\'olya Helix}
\author{}
\date{}

\begin{document}
\maketitle

\begin{abstract}
We introduce two space curves associated with the Riemann zeta function.  
The \emph{conical helix} interpolates the terms of the Dirichlet series of \(\zeta(s)\), each term expressed as a polar coordinate and placed one unit apart on a vertical axis.  
The \emph{Hilbert--P\'olya helix} (or Riemann helix) is a curve on the unit cylinder built from the non‑trivial zeros; its torsion equals a prime‑supported distribution.  
Together they provide a geometric framework for studying the relationship between the primes and the zeros.
\end{abstract}

\section{The Conical Helix of the Dirichlet Series}
For a fixed complex number \(s = \sigma + i t\) with \(\sigma > 1\), the Dirichlet series
\[
\zeta(s) = \sum_{n=1}^{\infty} \frac{1}{n^{s}}
\]
can be written term‑by‑term in polar form:
\[
\frac{1}{n^{s}} = \frac{1}{n^{\sigma}}\,e^{-i t \log n}
      = r_n \, e^{i \theta_n},
\qquad
r_n = n^{-\sigma},\quad \theta_n = -t \log n .
\]
We embed each term into \(\mathbb{R}^{3}\) by placing its polar coordinates in the horizontal plane and using the integer index as the vertical coordinate:
\[
P_n = \bigl( r_n \cos \theta_n,\; r_n \sin \theta_n,\; n \bigr),\qquad n = 1,2,3,\dots
\]

The points \(P_n\) lie at heights \(z = n\), with radii decaying like \(n^{-\sigma}\).  
A smooth curve that interpolates all the \(P_n\) is obtained by letting the index run over all real numbers \(x \ge 1\):
\[
\boxed{\;
C_{\text{Dir}}(x) \;=\; \bigl( x^{-\sigma} \cos(t \log x),\; 
                         x^{-\sigma} \sin(t \log x),\; 
                         x \bigr),\qquad x \ge 1 .
\;}
\]
This is a \emph{conical helix}:
\begin{itemize}
\item The distance from the \(z\)-axis is \(r(x) = x^{-\sigma}\).  For \(\sigma > 0\) the curve spirals inward as it ascends.
\item The angular coordinate \(\theta(x) = t \log x\) grows logarithmically, so the helix winds infinitely often as \(x \to \infty\).
\item The vertical spacing between successive integers is exactly \(1\); the curve passes through every \(P_n\).
\end{itemize}

When \(s\) is on the critical line \(\sigma = \frac12\), the radius decays as \(x^{-1/2}\); this is precisely the normalisation that appears in the prime‑sum side of the explicit formula.  
The partial sums \(\sum_{n=1}^{N} n^{-s}\) correspond to the polygonal path joining \(P_1, P_2, \dots, P_N\).

\section{The Hilbert--P\'olya Helix (Riemann Helix)}
Let the non‑trivial zeros of \(\zeta(s)\) with positive imaginary part be
\[
\rho_n = \beta_n + i \gamma_n,\qquad 0 < \gamma_1 < \gamma_2 < \cdots .
\]
Choose a smooth, strictly increasing function \(\phi: [\gamma_1,\infty) \to \mathbb{R}\) such that
\[
\phi(\gamma_n) = 2\pi n \qquad (n = 1,2,\dots).
\]
Such a phase function exists by monotone interpolation.  The \emph{Hilbert--P\'olya helix} (or unit‑radius Riemann helix) is the space curve
\[
\boxed{\;
C_{\text{HP}}(t) \;=\; \bigl( \cos \phi(t),\; \sin \phi(t),\; t \bigr),\qquad t \ge \gamma_1 .
\;}
\]
By construction, every zero is geometrically marked on the same generator of the cylinder:
\[
C_{\text{HP}}(\gamma_n) = (1,0,\gamma_n).
\]
The torsion of this curve is given by the Frenet–Serret formula
\[
\tau_{\text{HP}}(t) = 
\frac{\phi'(t)\bigl[ \phi'(t)^4 - \phi'(t)\phi'''(t) + 3\phi''(t)^2 \bigr]}
     {(\phi'(t)^2 + 1)\bigl( \phi'(t)^4 + \phi''(t)^2 + \phi'(t)^6 \bigr)} .
\]

Using the unconditional Weil explicit formula, one can show that, in the sense of distributions,
\[
\tau_{\text{HP}}(t) \;=\; \sum_{p,m} \frac{\log p}{p^{m/2}}\;\delta(t - m \log p) \;+\; \tau_{\text{arch}}(t),
\]
where the sum runs over all prime powers and \(\tau_{\text{arch}}\) is a smooth function coming from the archimedean Gamma factor.  
Thus the torsion of the Hilbert--P\'olya helix is exactly the prime‑supported distribution that appears on the right‑hand side of the explicit formula.

\section{Relation to the Riemann Hypothesis}
If one additionally encodes the real parts \(\beta_n\) by allowing the radius to vary — i.e. defining
\[
C_{r}(t) = \bigl( r(t)\cos\phi(t),\; r(t)\sin\phi(t),\; t \bigr),\qquad r(\gamma_n)=\beta_n,
\]
then the torsion \(\tau_{r}(t)\) contains additional singularities proportional to \(r'(t)\) and \(r''(t)\) at the prime logarithms.  The condition \(\tau_{r} = \tau_{\text{HP}}\) (equality of distributions) forces \(r'=0\) and \(r''=0\) on a dense set, hence \(r\) constant.  Symmetry from the functional equation then forces \(r \equiv \frac12\), which is the Riemann Hypothesis.

\section{Conclusion}
The conical helix provides a geometric encoding of the Dirichlet series terms; the Hilbert--P\'olya helix encodes the zeros.  The interplay between these two curves — one built from the Euler product data, the other from the zero data — mirrors the duality expressed by the explicit formula and offers a purely geometric formulation of the Riemann Hypothesis.

\end{document}